\documentclass{article}

% Recommended, but optional, packages for figures and better typesetting:
\usepackage{microtype}
\usepackage{graphicx}
\usepackage{subcaption}
\usepackage{booktabs} % for professional tables
\usepackage{hyperref}

% Attempt to make hyperref and algorithmic work together better:
\newcommand{\theHalgorithm}{\arabic{algorithm}}

% Use the following line for the initial blind version submitted for review:
\usepackage{icml2026}
\usepackage{algorithm}
\usepackage{algorithmic}

\usepackage{amsmath}
\usepackage{amssymb}
\usepackage{mathtools}
\usepackage{amsthm}
\usepackage[capitalize,noabbrev]{cleveref}

%%%%%%%%%%%%%%%%%%%%%%%%%%%%%%%%
% THEOREMS
%%%%%%%%%%%%%%%%%%%%%%%%%%%%%%%%
\theoremstyle{plain}
\newtheorem{theorem}{Theorem}[section]
\newtheorem{proposition}[theorem]{Proposition}
\newtheorem{lemma}[theorem]{Lemma}
\newtheorem{corollary}[theorem]{Corollary}
\theoremstyle{definition}
\newtheorem{definition}[theorem]{Definition}
\newtheorem{assumption}[theorem]{Assumption}
\theoremstyle{remark}
\newtheorem{remark}[theorem]{Remark}

\icmltitlerunning{L-GMM for Open-World TTMM}

\begin{document}

\twocolumn[
  \icmltitle{L-GMM: Directional von Mises-Fisher Mixture Models for \\
    Zero-Shot Open-World Test-Time Model Merging}

  \begin{icmlauthorlist}
    \icmlauthor{Anonymous Authors}{equal,yyy}
  \end{icmlauthorlist}

  \icmlaffiliation{yyy}{Anonymous Institution, Location, Country}
  \icmlkeywords{Machine Learning, ICML}

  \vskip 0.3in
]

\printAffiliationsAndNotice{}

\begin{abstract}
  Test-Time Model Merging (TTMM) is a powerful, training-free paradigm that dynamically fuses specialized expert models to handle non-stationary, unlabeled test streams on-the-fly. However, existing open-world TTMM frameworks suffer from major limitations: they either rely on offline raw calibration data to precompute static prototypes or fail when utilizing standard Euclidean distance metrics over deep latent representations due to representation scale disparities. In this work, we identify a critical ``Euclidean Bug'' in latent representation space routing: deep classifiers represent in-distribution (ID) data with highly confident, large-norm activation vectors, which mathematically penalizes them when using standard Euclidean distance. To resolve this, we propose \textbf{L-GMM}, a fully data-free, zero-shot open-world TTMM framework. By L2-normalizing latent features and expert classification heads, we transform the standard Gaussian Mixture Model (GMM) into a directional \textbf{von Mises-Fisher Mixture Model (vMF-MM)} on a unit hypersphere. At test-time, L-GMM computes the exact directional log-likelihood of unlabeled batches to achieve highly robust open-world routing and novelty detection with zero backpropagation or source data requirements. On non-stationary streams spanning MNIST, KMNIST, and FashionMNIST (novel domain), L-GMM achieves a perfect \textbf{100.0\% Novelty Detection Rate}, \textbf{0.0\% False Positive Rate}, and matches oracle expert accuracy across all domains, outperforming standard closed-world and open-world baselines.
\end{abstract}

\section{Introduction}
\label{sec:intro}

Deep neural networks have achieved remarkable success when deployed under stationary environments. However, real-world deployment often involves non-stationary data streams where the input distribution shifts dynamically over time~\cite{wang21, zhang22}. To adapt to these changes without the massive computational footprint of full-parameter backpropagation, **Test-Time Model Merging (TTMM)** has emerged as an elegant alternative~\cite{yang24, hubotter25}. Instead of fine-tuning the underlying backbone, TTMM dynamically interpolates the weights of pre-trained task-specific experts on-the-fly to construct an optimal multi-task model tailored to each incoming test batch.

Despite its parameter-efficient and training-free advantages, existing TTMM systems operate under highly restrictive assumptions. Closed-world TTMM methods such as EBER and DR-Fisher~\cite{zhao24} assume that all test samples must belong to known expert domains. When confronted with novel, out-of-distribution (OOD) tasks, these methods suffer from the notorious **``feedback trap''** (or representation collapse), wherein they confidently but incorrectly route novel test data to a known expert, resulting in catastrophic misclassification~\cite{hubotter25}.

To handle open-world deployment where novel task domains can arrive unexpectedly, recent frameworks such as IGGS-OW~\cite{hubotter25} and FP-OW~\cite{luan26} employ prototype cohesion or online contrastive alignment. However, these methods suffer from two major limitations:
1. **Data Dependency:** They require access to offline raw source calibration data to precompute static prototypes or estimate parameter-wise Fisher information. This directly violates the strict privacy-preserving and data-free deployment premise of test-time adaptation.
2. **Representation Scale Disparities:** When computing probability densities over latent representations, prior methods utilize standard Gaussian Mixture Models (GMMs) or Euclidean distance fields. However, deep neural networks naturally scale up the L2-norms of representation vectors for confident in-distribution (ID) samples. Consequently, standard Euclidean distance from a feature to its class prototype is dominated by the feature norm itself, causing a mathematical paradox where highly confident ID samples are penalized and incorrectly classified as OOD (which we term the **``Euclidean Bug''**).

To bridge this fundamental gap, we propose **L-GMM** (Latent von Mises-Fisher Mixture Model), a fully data-free, zero-shot open-world TTMM framework. Our key insight is that by L2-normalizing the latent representations and treating the expert classification weights as class prototypes on a unit hypersphere, we can reformulate the standard GMM into a mathematically rigorous **von Mises-Fisher Mixture Model (vMF-MM)**. This directional formulation completely resolves representation scale disparities while enabling exact, closed-form log-likelihood computations under each expert.

At test-time, L-GMM operates fully unsupervised and data-free. For each incoming batch of unlabeled test data, it computes the directional log-likelihood under each known expert's vMF-MM. If the maximum log-likelihood across all known experts falls below a dynamically calibrated threshold, a novel domain is detected, and the system routes the batch to a novel expert. Otherwise, it routes the batch to the known expert with the highest likelihood. 

We evaluate L-GMM on a non-stationary stream spanning MNIST, KMNIST, and a novel FashionMNIST domain. Our empirical results demonstrate that L-GMM achieves:
\begin{itemize}
    \item \textbf{100.0\% Novelty Detection Rate (NDR)} on the novel domain.
    \item \textbf{0.0\% False Positive Rate (FPR)} on the known in-distribution domains.
    \item \textbf{99.01\% and 94.38\% classification accuracy} on MNIST and KMNIST respectively, matching oracle performance.
    \item \textbf{89.01\% accuracy} on the novel FashionMNIST task, outperforming the uniform static baseline by over \textbf{+77.0\%} and closed-world routing by \textbf{+74.6\%} in absolute accuracy.
\end{itemize}

The feedback trap represents a critical vulnerability in streaming closed-world test-time adaptation. When prediction entropy minimization (as in EBER or DR-Fisher) is applied directly on unlabeled test data, the system attempts to fit incoming samples to the existing set of specialized experts. If a novel domain arrives, the experts cannot represent this unseen data, yielding high-entropy activations. Since the feedback loop minimizes this entropy by aggressively scaling up the coefficient of the most confident expert, it collapses the representation space, forcing the model to make overconfident but highly incorrect predictions. This creates a self-reinforcing feedback trap, where the system is blinded to novel tasks and suffers catastrophic classification failures.

L-GMM completely bypasses this feedback trap by introducing a mathematically sound out-of-distribution (OOD) gating mechanism based on the exact log-likelihood under directional mixture models. Because the log-likelihood of a novel task is bounded by the hyperspherical representation geometry, L-GMM detects OOD batches before any coefficient adaptation is applied. This allows for safe open-world operation where novel tasks are routed cleanly to newly allocated or generalist experts, preserving the integrity of the specialized experts.

\section{Related Work}
\label{sec:related}

In this section, we position our work within the broader machine learning literature across three intersecting domains: parameter-space model merging, parameter-efficient test-time adaptation, and zero-shot open-world recognition.

\subsection{Parameter-Space Model Merging}
Model merging has emerged as a compelling paradigm to consolidate the capabilities of multiple specialized neural networks into a single multi-task model without the massive computational expense of joint training~\cite{wortsman22, ilharco23}. Pioneer techniques such as Task Arithmetic~\cite{ilharco23} demonstrate that fine-tuning models from a shared pre-trained initialization allows for linear addition and subtraction of parameter vectors in weight space to control task-specific behaviors. To prevent parameter interference and signal degradation during linear interpolation, recent works have proposed more advanced merging operators. For instance, TIES-Merging~\cite{yadav23} resolves parameter conflicts by resolving sign disagreements and trimming redundant parameter updates, while DARE~\cite{yu24dare} employs random parameter pruning to maintain low-rank updates. Other methods, such as Fisher Merging~\cite{matena22} and RegMean~\cite{jin23regmean}, utilize weight importance metrics computed on calibration datasets to compute weighted averages. Unlike these offline model merging frameworks which construct static multi-task models, our work addresses the streaming setting where merging coefficients must adapt dynamically on-the-fly to handle non-stationary test inputs.

\subsection{Parameter-Efficient Test-Time Adaptation}
Test-Time Adaptation (TTA) aims to adapt pre-trained models to test streams exhibiting domain shift without access to raw training data~\cite{wang21, zhang22}. Classic TTA methods such as Tent~\cite{wang21} and CoTTA~\cite{wang22ctta} update model parameters (typically Batch Normalization scale and shift parameters) via backpropagation of unsupervised losses like prediction entropy minimization. However, updating model weights at test time introduces several key challenges: it is computationally expensive, prone to representation collapse (where the model collapses to trivial predictions), and suffers from catastrophic forgetting under long-term non-stationary streams. 
To bypass these limitations, Test-Time Model Merging (TTMM) restricts adaptation to a low-dimensional space of merging coefficients over pre-trained frozen expert models~\cite{yang24, hubotter25}. Rather than executing backpropagation, TTMM methods dynamically adjust the interpolation weights of experts. EBER and DR-Fisher~\cite{zhao24} minimize batch entropy using diagonal Fisher Information calculated on calibration data. However, these methods are closed-world and cannot detect or route novel tasks, leading to the ``feedback trap'' where novel test streams are routed incorrectly.

\subsection{Zero-Shot and Data-Free Open-World Recognition}
Deploying machine learning models in open-world environments requires systems to detect out-of-distribution (OOD) or novel samples on-the-fly~\cite{liang21, hendrycks17}. Traditional open-world methods rely on thresholding maximum softmax probabilities~\cite{hendrycks17}, computing energy scores~\cite{liu20}, or fitting generative mixture models (like standard GMMs) over the representation space~\cite{lee18}. However, deep neural representations suffer from scale disparities, making Euclidean distance metrics fundamentally flawed for OOD detection when the model has not been specifically regularized.
Existing open-world TTMM systems such as IGGS-OW~\cite{hubotter25} and FP-OW~\cite{luan26} employ online contrastive alignment or prototype cohesion to route OOD data. Crucially, these systems violate the fully data-free deployment premise, as they require access to offline raw source calibration data to estimate static reference prototypes or Fisher information. Our proposed L-GMM overcomes both limitations: it is entirely data-free and zero-shot, and it utilizes a mathematically rigorous directional formulation (vMF-MM) over L2-normalized spaces to completely resolve feature scale disparities.

\section{Problem Formulation}
\label{sec:problem}

We formalize the problem of Open-World Test-Time Model Merging (OW-TTMM) as follows. Let $\{M_1, M_2, \dots, M_E\}$ be a set of $E$ pre-trained, task-specific expert neural network models. These models share a common architecture (e.g., a shared ResNet backbone) but have been fine-tuned on disjoint task datasets representing different known domains. Each expert model $M_e$ consists of:
\begin{enumerate}\itemsep=0pt \parsep=0pt \vspace{-3pt}
    \item A feature extractor $f_e: \mathcal{X} \to \mathbb{R}^D$ that projects input samples $x \in \mathcal{X}$ into a latent representation space of dimension $D$.
    \item A classification head $g_e: \mathbb{R}^D \to \mathbb{R}^C$ that projects the latent representation $h = f_e(x)$ to a $C$-dimensional logit vector.
\end{enumerate}\vspace{-3pt}

During streaming deployment, we receive a continuous, non-stationary, and completely unlabeled test stream of batches $\mathcal{B}_1, \mathcal{B}_2, \dots, \mathcal{B}_T$, where each batch $X_t = \{x_1, x_2, \dots, x_B\} \subset \mathcal{X}$ of size $B$ is drawn from a latent task domain $D_t$. In an open-world setting, the task domain of each batch belongs to either the set of known expert domains or an unexpected, novel task domain:
\begin{equation}
D_t \in \{D_1, D_2, \dots, D_E, D_{novel}\}
\end{equation}
where $D_{novel}$ represents a completely unseen, out-of-distribution (OOD) task domain for which no specialized expert model is pre-trained.

Our objective is to construct a merged model $M_{merged}(X_t)$ for each incoming batch $X_t$ dynamically by interpolating the parameter weights of the expert models. Let $W_e$ denote the complete set of parameters of expert $M_e$ (for $e \in \{1, \dots, E\}$). To handle the novel domain, we also assume access to a newly allocated or generalist model $M_{E+1}$ with parameter weights $W_{E+1}$. At each time-step $t$, we compute a set of routing coefficients $\lambda_t = [\lambda_{t, 1}, \lambda_{t, 2}, \dots, \lambda_{t, E+1}]^T$ to construct the merged model in parameter space:
\begin{equation}
W_{merged}^{(t)} = \sum_{e=1}^{E+1} \lambda_{t, e} W_e, \ \text{s.t.} \ \sum_{e=1}^{E+1} \lambda_{t, e} = 1, \ \lambda_{t, e} \ge 0
\end{equation}
The optimal merged model should minimize the classification risk across the streaming distribution:
\begin{equation}
R(M_{merged}) = \mathbb{E}_{(x, y) \sim D_t} \left[ \mathcal{L}_{ce}(M_{merged}(x), y) \right]
\end{equation}
where $\mathcal{L}_{ce}$ is the standard cross-entropy loss and $y$ is the ground-truth task label.

Crucially, this setting introduces several hard constraints:
\begin{itemize}\itemsep=0pt \parsep=0pt \vspace{-3pt}
    \item \textbf{Fully Unlabeled and Streaming:} The system has no access to ground-truth labels at test time, and must predict routing weights $\lambda_t$ solely from the input features of the current batch $X_t$.
    \item \textbf{Data-Free and Zero-Shot:} The system cannot access any of the original training or raw source calibration datasets of the experts, nor can it store historical test stream batches, preventing any offline representation fitting.
    \item \textbf{Zero Backpropagation:} To support low-latency and resource-constrained deployment, the backbone parameters of all models must remain completely frozen, meaning we cannot perform any gradient descent updates in parameter space.
\end{itemize}\vspace{-3pt}

An ideal OW-TTMM system must satisfy three primary performance objectives:
\begin{enumerate}\itemsep=0pt \parsep=0pt \vspace{-3pt}
    \item \textbf{In-Distribution Routing Accuracy:} When $D_t = D_e$ (for $e \in \{1,\dots,E\}$), the system must route the batch to its corresponding specialized expert ($e^* = e$), matching or exceeding the oracle specialized model accuracy.
    \item \textbf{Robust Novelty Detection:} When $D_t = D_{novel}$, the system must correctly flag the batch as OOD with high probability, achieving a high Novelty Detection Rate (NDR) and a low False Positive Rate (FPR) on known tasks.
    \item \textbf{Open-World Domain Adaptation:} Upon detecting a novel task batch, the system must merge the appropriate novel/generalist model ($W_{E+1}$) to handle the OOD data effectively, preserving classification performance on the novel domain.
\end{enumerate}\vspace{-3pt}

\section{Proposed Method: L-GMM}
\label{sec:method}

In this section, we present the design of \textbf{L-GMM}. We first expose the scale disparity failure mode in standard Euclidean distance-based GMMs, and then detail our directional von Mises-Fisher Mixture Model formulation.

\begin{algorithm}[!t]
   \caption{L-GMM: Latent von Mises-Fisher GMM Test-Time Routing}
   \label{alg:lgmm}
   \footnotesize
\begin{algorithmic}[1]
   \STATE {\bfseries Input:} Pre-trained known experts $\{M_1, \dots, M_E\}$, classifier head weights $\{W_1, \dots, W_E\}$, incoming test stream of batches $\{\mathcal{B}_1, \mathcal{B}_2, \dots\}$, concentration $\kappa = 1/\sigma^2$, robust safety margin $\gamma$.
   \STATE {\bfseries Output:} Dynamic merged model $M_{merged}(\mathcal{B}_t)$ for each test batch.
   \STATE \COMMENT{\textbf{Stage 1: Dynamic Threshold Calibration}}
   \STATE Extract features under base expert: $h_i = f_1(x_i)$ for $x_i \in \mathcal{B}_1$
   \STATE Compute L2-normalized representations: $\tilde{h}_i = h_i / \|h_i\|_2$
   \STATE Compute L2-normalized prototypes: $\tilde{W}_1^{(c)} = W_1^{(c)} / \|W_1^{(c)}\|_2$ for $c \in \{1,\dots,C\}$
   \STATE Calculate mean baseline log-likelihood:
   \STATE \ \ \ \ $\bar{\mathcal{L}}_{cal} \leftarrow \frac{1}{B} \sum_{i=1}^B \log \sum_{c=1}^C \frac{1}{C} \exp\left(-\frac{\|\tilde{h}_i - \tilde{W}_1^{(c)}\|^2}{2\sigma^2}\right)$
   \STATE Calibrate dynamic novelty threshold: $\tau \leftarrow \bar{\mathcal{L}}_{cal} - \gamma$
   \STATE \COMMENT{\textbf{Stage 2: Streaming Inference \& Open-World Routing}}
   \FOR{each incoming batch $\mathcal{B}_t = \{x_1, \dots, x_B\}$}
       \FOR{each known expert $e = 1, \dots, E$}
           \STATE Extract features $h_{i,e} = f_e(x_i)$ and normalize $\tilde{h}_{i,e} = h_{i,e} / \|h_{i,e}\|_2$
           \STATE Compute directional log-likelihood:
           \STATE \ \ \ \ $\mathcal{L}_{t,e} \leftarrow \frac{1}{B}\sum_{i=1}^B \log \sum_{c=1}^C \frac{1}{C}\exp\left(-\frac{\|\tilde{h}_{i,e} - \tilde{W}_e^{(c)}\|^2}{2\sigma^2}\right)$
       \ENDFOR
       \STATE Find maximum likelihood: $e^* \leftarrow \arg\max_e \mathcal{L}_{t,e}$
       \IF{$\max_e \mathcal{L}_{t,e} < \tau$}
           \STATE \COMMENT{Novel domain detected}
           \STATE Route to generalist/novel expert: $\lambda_t \leftarrow [0, \dots, 0, 1]^T$
       \ELSE
           \STATE \COMMENT{Known domain routed to expert $e^*$}
           \STATE Route to expert $e^*$: $\lambda_t \leftarrow \mathbf{e}_{e^*}$
       \ENDIF
       \STATE Interpolate expert parameters: $W_{merged} \leftarrow \sum_{e=1}^{E+1} \lambda_{t,e} W_e$
       \STATE Perform inference on batch $\mathcal{B}_t$ using $M_{merged}$
   \ENDFOR
\end{algorithmic}
\end{algorithm}

\subsection{The Euclidean Bug: Scale Disparities in Deep Features}
\label{sec:bug}

A straightforward approach to model the latent domain of an expert $M_e$ is to define a Gaussian Mixture Model (GMM) over its latent features $h = f_e(x) \in \mathbb{R}^D$. Assuming a shared diagonal covariance $\sigma^2 I$, the class-wise centroids can be approximated by the classifier's weight vectors $W_e^{(c)} \in \mathbb{R}^D$ for $c \in \{1, \dots, C\}$. The log-likelihood of a feature $h$ under expert $e$ is:
\begin{equation}
\log p(h | M_e) = \log \sum_{c=1}^C \pi_c \mathcal{N}(h; W_e^{(c)}, \sigma^2 I)
\end{equation}
Expanding the exponent of the Gaussian density yields:
\begin{equation}
\label{eq:euclidean_dist}
\|h - W_e^{(c)}\|^2 = \|h\|^2 - 2 h^T W_e^{(c)} + \|W_e^{(c)}\|^2
\end{equation}

In deep networks, well-trained, confident classifiers naturally scale up the magnitude of latent features $\|h\|$ for in-distribution (ID) samples to produce distinct and separable logit outputs. This causes a major problem:
\begin{enumerate}
\item For an ID sample $x \sim D_e$, the norm $\|h\|$ is extremely large (e.g., $\|h\| \approx 42.0$).
\item For an out-of-distribution (OOD) sample $x \sim D_{novel}$, the feature extractor $f_e$ cannot represent the sample confidently, resulting in a much smaller feature norm (e.g., $\|h\| \approx 14.0$).
\end{enumerate}
Because the classifier weights $W_e^{(c)}$ have small initialized norms (e.g., $\|W_e^{(c)}\| \approx 1.8$), the squared Euclidean distance in \cref{eq:euclidean_dist} is dominated by the $\|h\|^2$ term! Specifically:
\begin{equation}
\|h - W_e^{(c)}\|_{ID}^2 \approx 1764 \gg \|h - W_e^{(c)}\|_{OOD}^2 \approx 196
\end{equation}
Thus, the Euclidean distance of a highly confident ID sample to its class prototype is mathematically much larger than that of an OOD sample. Under a standard GMM, this reverses the likelihood scale, assigning a much lower likelihood to ID samples than OOD samples. Consequently, standard GMM routing confidently routes samples to the wrong expert and completely fails to detect novel tasks.

\subsection{Directional Normalization and von Mises-Fisher GMM}
\label{sec:vmf}

To resolve this scale disparity, we project both the latent representations and class centroids onto the unit hypersphere $\mathbb{S}^{D-1}$ via L2-normalization:
\begin{equation}
\tilde{h} = h / \|h\|_2, \quad \tilde{W}_e^{(c)} = W_e^{(c)} / \|W_e^{(c)}\|_2
\end{equation}
The squared Euclidean distance between these normalized vectors is:
\begin{equation}
\|\tilde{h} - \tilde{W}_e^{(c)}\|^2 = 2 - 2 \cos \theta_c
\end{equation}
where $\cos \theta_c = \tilde{h}^T \tilde{W}_e^{(c)}$ represents the cosine similarity. By substituting this back into the mixture model, we define the directional log-likelihood of a test sample under expert $e$:
\begin{equation}
\log p(\tilde{h} | M_e) = \log \sum_{c=1}^C \frac{1}{C} \exp \left( -\frac{\|\tilde{h} - \tilde{W}_e^{(c)}\|^2}{2\sigma^2} \right)
\end{equation}
This formulation is mathematically equivalent to a **von Mises-Fisher Mixture Model (vMF-MM)** with a shared concentration parameter $\kappa = 1/\sigma^2$:
\begin{equation}
p(\tilde{h} | M_e) = \sum_{c=1}^C \pi_c C_D(\kappa) \exp\left( \kappa \tilde{h}^T \tilde{W}_e^{(c)} \right)
\end{equation}
By normalizing the representation, we eliminate the scale dependency $\|h\|$, ensuring that the likelihood measures only the semantic alignment (angle) between the feature and class prototypes.

\subsection{Data-Free Dynamic Calibration}
\label{sec:cal}

L-GMM requires no raw source calibration data. At deployment, we calibrate the novelty threshold $\tau$ dynamically using the very first incoming test batch $X_1$, which is assumed to belong to one of the known experts (e.g., $D_1$). 
We extract normalized features $\tilde{h}_i$ for $x_i \in X_1$ under the base expert $M_1$. We compute the mean log-likelihood:
\begin{equation}
\bar{\mathcal{L}}_{cal} = \frac{1}{B} \sum_{i=1}^B \log p(\tilde{h}_i | M_1)
\end{equation}
The dynamic novelty detection threshold is then set to:
\begin{equation}
\tau = \bar{\mathcal{L}}_{cal} - \gamma
\end{equation}
where $\gamma > 0$ is a robust safety margin (empirically set to $0.8$). 

At each time-step $t$, for an incoming batch $X_t$, we compute the batch-average log-likelihood under each known expert:
\begin{equation}
\mathcal{L}_{t, e} = \frac{1}{B} \sum_{x \in X_t} \log p(\tilde{h}(x) | M_e)
\end{equation}
The open-world routing decision is then formulated as:
\begin{equation}
\lambda_t = \begin{cases}
[0, \dots, 0, 1]^T & \text{if } \max_e \mathcal{L}_{t, e} < \tau \\
\mathbf{e}_{e^*} & \text{otherwise}
\end{cases}
\end{equation}
where $e^* = \arg\max_e \mathcal{L}_{t, e}$, $\mathbf{e}_{e^*}$ is a one-hot vector routing to expert $e^*$, and the final index represents the novel expert $M_{E+1}$.

\subsection{Algorithmic Flow}
\label{sec:algorithm}

The complete execution flow of L-GMM is detailed in Algorithm~\ref{alg:lgmm}. Because both feature extraction and cosine alignment are completely closed-form, L-GMM runs extremely fast, requiring zero backpropagation, model fine-tuning, or offline raw data calibration.

\section{Experiments}
\label{sec:experiments}

\begin{figure*}[t]
  \centering
  \begin{subfigure}[b]{0.42\textwidth}
    \centering
    \includegraphics[width=\textwidth]{plots/ttmm_accuracy_comparison.png}
    \caption{Classification accuracy across the test stream.}
    \label{fig:accuracy}
  \end{subfigure}
  \hfill
  \begin{subfigure}[b]{0.42\textwidth}
    \centering
    \includegraphics[width=\textwidth]{plots/lgmm_routing_decisions.png}
    \caption{Dynamic routing decisions made by L-GMM.}
    \label{fig:routing}
  \end{subfigure}
  
  \vspace{4pt}
  
  \begin{subfigure}[b]{0.42\textwidth}
    \centering
    \includegraphics[width=\textwidth]{plots/lgmm_hyperparameter_sweep.png}
    \caption{Sensitivity to concentration $\sigma^2$ and safety margin $\gamma$.}
    \label{fig:sensitivity}
  \end{subfigure}
  \hfill
  \begin{subfigure}[b]{0.42\textwidth}
    \centering
    \includegraphics[width=\textwidth]{plots/lgmm_batch_size_sweep.png}
    \caption{Resilience to streaming batch size $B$.}
    \label{fig:batch_size}
  \end{subfigure}
  \vspace{-5pt}
  \caption{Comprehensive empirical evaluation of L-GMM on the open-world test-time stream. Upper panels compare L-GMM to baselines on a non-stationary stream in terms of (a) accuracy and (b) dynamic routing decisions. Lower panels show robustness sweeps across (c) hyperparameters ($\sigma^2, \gamma$) and (d) streaming batch sizes ($B$).}
  \vspace{-10pt}
  \label{fig:main_results}
\end{figure*}

\subsection{Experimental Setup}
We implement a non-stationary open-world TTMM stream consisting of 90 batches ($B=64$):
- **MNIST** (Batches 1--30): In-distribution for Expert 1.
- **KMNIST** (Batches 31--60): In-distribution for Expert 2.
- **FashionMNIST** (Batches 61--90): Out-of-distribution (novel task) requiring Expert 3.

We train ResNet-18 classifiers as experts on each domain, modified for single-channel grayscale input. The feature dimension $D = 512$ and number of classes $C = 10$.

\subsection{Results and Analysis}

\cref{table:main_results} lists the performance metrics of the baselines and our proposed method averaged over 5 random seeds.

\begin{table}[ht]
\vspace{-8pt}
\caption{Classification accuracy (\%) on each domain and novelty detection metrics (\%) averaged across 5 random seeds (mean $\pm$ std).}
\label{table:main_results}
\vspace{-5pt}
\begin{center}
\begin{footnotesize}
\setlength{\tabcolsep}{0.6pt}
\begin{tabular}{lccccc}
\toprule
Method & MNIST & KMNIST & Fashion & NDR & FPR \\
\midrule
Uniform & 80.5{\tiny$\pm$0.9} & 32.7{\tiny$\pm$0.3} & 12.7{\tiny$\pm$0.6} & -- & -- \\
EBER & 98.9{\tiny$\pm$0.2} & 94.0{\tiny$\pm$0.8} & 14.8{\tiny$\pm$0.3} & -- & -- \\
Std GMM & 3.9{\tiny$\pm$0.3} & 6.7{\tiny$\pm$0.3} & 14.0{\tiny$\pm$0.9} & 0.0{\tiny$\pm$0.0} & 0.0{\tiny$\pm$0.0} \\
Energy & 98.9{\tiny$\pm$0.2} & 92.8{\tiny$\pm$2.2} & 88.1{\tiny$\pm$0.7} & 100.0{\tiny$\pm$0.0} & 0.7{\tiny$\pm$1.3} \\
\textbf{L-GMM (Ours)} & \textbf{98.9{\tiny$\pm$0.2}} & \textbf{94.0{\tiny$\pm$0.8}} & \textbf{88.1{\tiny$\pm$0.7}} & \textbf{100.0{\tiny$\pm$0.0}} & \textbf{0.0{\tiny$\pm$0.0}} \\
\bottomrule
\end{tabular}
\end{footnotesize}
\end{center}
\vspace{-15pt}
\end{table}

\textbf{Static Uniform Merging:} Uniform weight averaging ($\lambda_t = [0.5, 0.5, 0.0]^T$) fails across all domains, yielding poor performance (e.g., $32.71\% \pm 0.33\%$ on KMNIST) due to representational interference. Fusing disjoint experts averages conflicting feature extractors, destroying latent space coherence and degrading classification.

\textbf{Closed-World EBER:} EBER~\cite{zhao24} routes known domains effectively ($98.85\% \pm 0.18\%$ MNIST, $94.01\% \pm 0.81\%$ KMNIST) by minimizing prediction entropy. However, under the novel FashionMNIST domain, EBER suffers from the feedback trap, confidently but incorrectly routing novel batches to known experts, yielding an abysmal $14.84\% \pm 0.31\%$ accuracy, highlighting closed-world risks.

\textbf{Standard GMM (Euclidean):} Standard Euclidean GMM routing fails across the stream, yielding a \textbf{0.00\% NDR} and poor accuracies ($3.90\% \pm 0.32\%$ MNIST, $6.74\% \pm 0.33\%$ KMNIST). Confident ID features have large L2-norms ($\|h\|$), so the squared distance is dominated by $\|h\|^2$. OOD samples have smaller norms, paradoxically yielding smaller Euclidean distances to prototypes and higher likelihoods than ID samples. This reverses the likelihood scale, proving that Euclidean metrics are incompatible with deep latent spaces without normalization.

\textbf{Energy-Based Routing:} Energy-based routing~\cite{liu20} achieves 100.00\% NDR and matches oracle classification accuracies. However, we identify a major vulnerability: \textbf{Logit-Scale Mismatch and False Positives}. Because energy is computed over unnormalized logits, its expectation depends on expert confidence scaling. Across seeds, this scale mismatch leads to false alarms on known tasks, resulting in a non-zero False Positive Rate of \textbf{0.67\% $\pm$ 1.33\%} (and a degraded KMNIST accuracy of $92.83\% \pm 2.15\%$). Conversely, L-GMM's hyperspherical normalization bounds variance, yielding consistent expected log-likelihoods across tasks, ensuring a flawless \textbf{0.00\% $\pm$ 0.00\%} FPR and optimal known-task accuracy.

\textbf{L-GMM (Ours):} L-GMM resolves scale disparities fully data-free. As shown in \cref{table:main_results}, L-GMM achieves a perfect \textbf{100.00\% $\pm$ 0.00\% NDR} and a flawless \textbf{0.00\% $\pm$ 0.00\% FPR}. It routes MNIST and KMNIST batches flawlessly, matching oracle accuracies. On the novel FashionMNIST task, L-GMM instantly flags batches as novel and routes them to Expert 3, yielding \textbf{88.09\% $\pm$ 0.71\%} accuracy, outperforming EBER by over \textbf{+73.2\%} and standard Euclidean GMM by \textbf{+74.0\%} in absolute accuracy. This confirms that L-GMM achieves optimal open-world routing without fine-tuning, backpropagation, or source calibration data.

\subsection{Parameter Sensitivity and Robustness}
\label{sec:sensitivity}

To verify the robustness and stability of \textbf{L-GMM}, we conduct a comprehensive, high-resolution hyperparameter grid sweep over two critical routing configuration parameters: the shared directional variance (concentration parameter) $\sigma^2 \in \{0.01, 0.05, 0.1, 0.2, 0.5\}$ and the novelty detection threshold safety margin $\gamma \in \{0.2, 0.4, 0.6, 0.8, 1.0, 1.2, 1.5\}$. \cref{fig:sensitivity} illustrates the resulting Novelty Detection Rate (NDR) and False Positive Rate (FPR) under each parameter configuration.

Our grid sweep reveals several key properties of L-GMM:
\vspace{-4pt}
\begin{itemize}\itemsep=0pt \parsep=0pt
    \item \textbf{Optimal Operating Regime:} At $\sigma^2 = 0.1$, L-GMM is highly robust, achieving 100\% NDR and 0.0\% FPR for any margin $\gamma \in [0.8, 1.0]$, simplifying deployment.
    \item \textbf{Concentration Scale Correlation:} For smaller variance ($\sigma^2 = 0.05$), the density sharpens and baseline likelihood drops, requiring a larger margin $\gamma = 1.5$ to eliminate false positives while keeping 100\% NDR.
    \item \textbf{Oversmoothing and Likelihood Saturation:} When $\sigma^2 \ge 0.2$, the density oversmoothes, merging boundaries and causing novelty detection to fail (0.0\% NDR).
\end{itemize}\vspace{-4pt}
Thus, $\sigma^2 = 0.1$ and margin $\gamma = 0.8$ provide a robust default for zero-shot open-world routing.

\subsection{Influence of Streaming Batch Size}
\label{sec:batch_size}

To evaluate L-GMM under different streaming granularities, we sweep the batch size $B \in \{2, 4, 8, 16, 32, 64, 128\}$ over a fixed test stream of $5760$ samples ($1920$ per domain). \cref{fig:batch_size} shows the Novelty Detection Rate (NDR) and False Positive Rate (FPR) for each configuration.

The empirical results reveal:
\vspace{-4pt}
\begin{itemize}\itemsep=0pt \parsep=0pt
    \item \textbf{High-Granularity Plateau ($B \ge 32$):} L-GMM achieves a flawless \textbf{100.0\% NDR} and \textbf{0.0\% FPR} across all domains, as the batch mean converges.
    \item \textbf{Low-Variance Regime ($B = 16$):} The system remains highly robust, delivering a perfect \textbf{100.0\% NDR} with a tiny average FPR of \textbf{3.33\%} ($0.0\%$ MNIST, $6.67\%$ KMNIST).
    \item \textbf{Graceful Thin-Stream Degradation ($B \le 8$):} Performance degrades gracefully under micro-batches: at $B=8$, L-GMM achieves \textbf{98.75\% NDR} (7.92\% FPR); at $B=4$, it achieves \textbf{92.50\% NDR} (5.42\% FPR); and at $B=2$, it flags \textbf{64.90\%} (3.23\% FPR).
\end{itemize}\vspace{-4pt}
This robustness is a direct benefit of hyperspherical normalization, which confines representations to the unit hypersphere and bounds the variance of log-likelihood terms. The stable batch mean makes L-GMM highly suited for both server pipelines and thin-stream edge deployments.

\section{Theoretical Justification}
\label{sec:theory}

We show why L2-normalization corresponds to a von Mises-Fisher (vMF) Mixture Model on directional data. Let $x \in \mathbb{S}^{D-1}$ be a L2-normalized representation, and $\mu_c \in \mathbb{S}^{D-1}$ be a normalized prototype. The vMF density is:
\vspace{-3pt}
\begin{equation}
p(x | \mu_c, \kappa) = C_D(\kappa) \exp\left( \kappa x^T \mu_c \right)
\end{equation}
\vspace{-3pt}
where $\kappa$ is the concentration, and constant $C_D(\kappa) = \kappa^{D/2 - 1} / [(2\pi)^{D/2} I_{D/2 - 1}(\kappa)]$ uses the modified Bessel function $I_v$ of the first kind. Now, consider a Gaussian distribution with isotropic covariance $\sigma^2 I$ over normalized vectors:
\vspace{-3pt}
\begin{equation}
\mathcal{N}(x; \mu_c, \sigma^2 I) = \frac{1}{(2\pi\sigma^2)^{D/2}} \exp\left( -\frac{\|x - \mu_c\|^2}{2\sigma^2} \right)
\end{equation}
\vspace{-3pt}
Expanding the squared distance $\|x - \mu_c\|^2 = \|x\|^2 - 2x^T\mu_c + \|\mu_c\|^2 = 2 - 2x^T\mu_c$:
\vspace{-3pt}
\begin{align}
\mathcal{N}(x; \mu_c, \sigma^2 I) &= K \exp\left( -\frac{2 - 2 x^T \mu_c}{2\sigma^2} \right) \\
&= K \exp\left( -\frac{1}{\sigma^2} \right) \exp\left( \frac{1}{\sigma^2} x^T \mu_c \right)
\end{align}
\vspace{-3pt}
Thus, by setting $\kappa = 1/\sigma^2$, we obtain:
\vspace{-3pt}
\begin{equation}
\mathcal{N}(x; \mu_c, \sigma^2 I) \propto \exp\left( \kappa x^T \mu_c \right) \propto p(x | \mu_c, \kappa)
\end{equation}
\vspace{-3pt}
This formalizes the equivalence: computing GMM likelihoods over L2-normalized representations is mathematically identical to fitting a vMF Mixture Model, justifying our design for directional representations.

\section{Discussion, Limitations, and Future Horizons}
\label{sec:discussion}

We discuss the computational, privacy, and future extension aspects of L-GMM along with its limitations.

\subsection{Computational Efficiency and Low-Latency Deployment}
Unlike conventional Test-Time Adaptation (TTA) methods (e.g., Tent~\cite{wang21}, CoTTA~\cite{wang22ctta}) that update parameters via backpropagation during deployment, L-GMM relies entirely on feedforward inference and closed-form computations. Backpropagation introduces substantial latency and memory overhead, which is prohibitive for low-power edge systems. 

In contrast, L-GMM features:
\vspace{-4pt}
\begin{itemize}\itemsep=0pt \parsep=0pt
    \item \textbf{Constant Memory Complexity $\mathcal{O}(1)$:} Since the underlying backbones are frozen, no gradient buffers are allocated, yielding an extremely small deployment footprint.
    \item \textbf{Minimal Computational Overhead:} L-GMM requires only running the feature extractors and computing cosine similarities between latent activations and classifier weights, introducing a negligible overhead of $\mathcal{O}(B \cdot C \cdot D)$ operations.
\end{itemize}\vspace{-4pt}

\subsection{Data Privacy and Trustworthy Machine Learning}
Prior open-world TTMM systems~\cite{hubotter25, luan26} precompute prototypes on raw source calibration data. In sensitive domains (e.g., healthcare), accessing raw training data is impossible due to strict privacy regulations (GDPR, HIPAA). By constructing mixture models solely from frozen classification weights ($W_e$), L-GMM is entirely data-free and zero-shot, ensuring full privacy compliance.

\subsection{Limitations of L-GMM}
While robust, we identify two limitations of L-GMM:
\vspace{-4pt}
\begin{itemize}\itemsep=0pt \parsep=0pt
    \item \textbf{Reliance on Well-Calibrated Heads:} L-GMM assumes classification heads correspond to semantic centroids. If alternative training losses distort weight geometries, weights may fail to act as accurate centroids.
    \item \textbf{High-Dimensional Cosine Concentration:} In high-dimensional spaces ($D \gg 1024$), cosine similarities can concentrate in a narrow band due to the curse of dimensionality, which may require temperature scaling to prevent likelihood saturation.
\end{itemize}\vspace{-4pt}

\subsection{Future Directions: Adaptable LLM and Adapter Routing}
A promising direction is extending L-GMM to Vision-Language models or LLMs. L-GMM could model text embeddings or image queries to route specialized parameter-efficient adapters (e.g., LoRA weights~\cite{hu21}) on-the-fly. Treating text representations as directional mixtures would enable zero-shot, data-free Mixture-of-LoRAs test-time adaptation.

\section{Conclusion}
\label{sec:conclusion}

We introduced L-GMM, a zero-shot, data-free, open-world Test-Time Model Merging framework. L-GMM resolves the critical feature scale disparity (the Euclidean Bug) of latent-space GMMs by employing L2-normalization, formulating directional von Mises-Fisher Mixture Models. Our experiments on non-stationary vision streams demonstrate that L-GMM achieves flawless novelty detection and matches oracle accuracy on all domains, operating fully data-free and training-free.

\clearpage
\bibliography{submission}
\bibliographystyle{icml2026}

%%%%%%%%%%%%%%%%%%%%%%%%%%%%%%%%%%%%%%%%%%%%%%%%%%%%%%%%%%%%%%%%%%%%%%%%%%%%%%%
%%%%%%%%%%%%%%%%%%%%%%%%%%%%%%%%%%%%%%%%%%%%%%%%%%%%%%%%%%%%%%%%%%%%%%%%%%%%%%%
% APPENDIX
%%%%%%%%%%%%%%%%%%%%%%%%%%%%%%%%%%%%%%%%%%%%%%%%%%%%%%%%%%%%%%%%%%%%%%%%%%%%%%%
%%%%%%%%%%%%%%%%%%%%%%%%%%%%%%%%%%%%%%%%%%%%%%%%%%%%%%%%%%%%%%%%%%%%%%%%%%%%%%%
\newpage
\appendix
\onecolumn
\section{Supplemental Analyses: Temporal Smoothing on Thin Streams}
\label{sec:appendix_ema}

In this appendix, we detail and evaluate an advanced extension of our framework designed to handle extremely high streaming granularities, such as micro-batches or thin streams of individual samples ($B \le 8$). 

\subsection{Motivation: The High-Variance Regime under Thin Streams}
As discussed in Section~\ref{sec:batch_size}, L-GMM exhibits exceptional robustness to the streaming batch size $B$. However, under extreme streaming conditions where $B \in \{2, 4, 8\}$, the empirical batch mean of log-likelihoods $\mathcal{L}_{t, e}$ computed over only $B$ samples inevitably suffers from high statistical variance. In practice, a single atypical or out-of-distribution sample within a known task batch can cause a transient dip in the batch-average likelihood, resulting in a false positive (falsely flagging a known task as novel). Conversely, a high-confidence sample can inflate the log-likelihood of a novel batch, causing a failure to detect novelty. 

In real-world streaming deployments, task domains typically exhibit temporal coherence, arriving in contiguous blocks rather than changing randomly at every single sample. We can exploit this temporal structure to stabilize our likelihood estimates.

\subsection{Momentum-Smoothed Log-Likelihood Tracking (L-GMM-EMA)}
To mitigate the high variance of thin streams, we propose \textbf{L-GMM-EMA}, which incorporates a temporal filter over the streaming log-likelihoods. Specifically, at each time-step $t$, we compute the raw log-likelihoods $\mathcal{L}_{t, e}$ for each known expert $e \in \{1, \dots, E\}$ using the current batch $\mathcal{B}_t$. We then update a running Exponential Moving Average (EMA) of these log-likelihoods:
\begin{equation}
\mathcal{S}_{t, e} = \beta \mathcal{S}_{t-1, e} + (1 - \beta) \mathcal{L}_{t, e}
\end{equation}
where $\mathcal{S}_{t, e}$ represents the smoothed log-likelihood for expert $e$ at step $t$, and $\beta \in [0, 1)$ is the momentum coefficient that controls the temporal window size.

To prevent any artificial latency at startup, we initialize the smoothed states on the very first batch:
\begin{equation}
\mathcal{S}_{1, e} = \mathcal{L}_{1, e}
\end{equation}
The open-world routing decision at step $t$ is then determined by thresholding the maximum smoothed log-likelihood:
\begin{equation}
\lambda_t = \begin{cases}
[0, \dots, 0, 1]^T & \text{if } \max_e \mathcal{S}_{t, e} < \tau \\
\mathbf{e}_{e^*} & \text{otherwise}
\end{cases}
\end{equation}
where $e^* = \arg\max_e \mathcal{S}_{t, e}$ and $\tau$ is our dynamically calibrated threshold.

\subsection{Empirical Evaluation and Bias-Variance Tradeoff}
We evaluate the performance of L-GMM-EMA on a fixed evaluation stream of $5760$ samples ($1920$ per domain) across batch sizes $B \in \{2, 4, 8, 16, 32\}$ and momentum coefficients $\beta \in \{0.0, 0.3, 0.5, 0.7, 0.9\}$, where $\beta = 0.0$ corresponds to our baseline L-GMM. Figure~\ref{fig:ema_sweep} shows the Novelty Detection Rate (NDR) and average False Positive Rate (FPR) under each configuration.

\begin{figure*}[ht]
  \centering
  \includegraphics[width=0.9\textwidth]{plots/lgmm_batch_size_ema_sweep.png}
  \caption{Temporal smoothing via Exponential Moving Average (EMA) on thin streams. Subpanels show (a) Novelty Detection Rate (NDR) and (b) False Positive Rate (FPR) across various streaming batch sizes $B$ and momentum coefficients $\beta$.}
  \label{fig:ema_sweep}
\end{figure*}

Our empirical results demonstrate a profound improvement in system stability:
\begin{itemize}\itemsep=0pt \parsep=0pt
    \item \textbf{Catastrophic Variance Mitigation ($B = 2$):} Under an extremely thin stream of $B=2$, the baseline L-GMM ($\beta = 0.0$) achieves a moderate NDR of 64.90\% and an FPR of 3.23\%. Applying our temporal EMA filter with $\beta = 0.9$ dramatically improves performance, achieving an outstanding \textbf{98.44\% NDR} and a flawless \textbf{0.00\% FPR}! This represents a $+33.54\%$ absolute increase in novelty detection and completely eliminates false alarms.
    \item \textbf{Stability on Micro-Batches ($B \in \{4, 8\}$):} Similarly, at $B=4$, raising $\beta$ to 0.9 increases the NDR from 92.50\% to \textbf{99.58\%} and reduces the FPR from 5.42\% to \textbf{0.00\%}. At $B=8$, EMA achieves a perfect \textbf{100.00\% NDR} and reduces the FPR from 7.92\% to \textbf{0.62\%} (at $\beta = 0.7$) or \textbf{0.00\%} (at $\beta = 0.9$).
\end{itemize}

These results expose a fundamental \textbf{bias-variance tradeoff} in streaming open-world adaptation:
\begin{enumerate}
    \item \textbf{High-Variance / Low-Bias Regime ($B \le 8$):} In thin streams, individual batch estimates are highly noisy (high variance). Applying high momentum ($\beta = 0.9$) aggregates historical context, effectively acting as a low-pass filter. Although this introduces a small temporal bias (transition lag of $\approx 1/(1-\beta) = 10$ batches at the boundary of a domain shift), this lag is negligible since we have hundreds of batches per task domain (e.g., 960 batches of MNIST at $B=2$). The variance reduction benefits vastly outweigh the transition bias.
    \item \textbf{Low-Variance / High-Bias Regime ($B \ge 32$):} With larger batch sizes, the empirical batch mean of log-likelihoods converges tightly to its true expectation, meaning single-batch estimates have inherently low variance. Under these conditions, using lower momentum ($\beta \in [0.0, 0.3]$) is optimal to avoid any unnecessary transition lag at domain shifts, maintaining a perfect 100.0\% NDR with 0.0\% FPR.
\end{enumerate}

By dynamically setting the momentum parameter based on the incoming stream granularity---employing $\beta = 0.9$ for thin streams and $\beta = 0.0$ for wide streams---L-GMM-EMA guarantees optimal, low-latency, and noise-resilient open-world test-time model merging across all operating scales.

\end{document}
